% GENERATED FILE. Edit canonical lessons, then run npm run generate:books.
% canonical-source-hash: fnv1a64:41c86f7063f6f8c0
% canonical-lessons: JA-C391-oudanhodou, JA-C391-hokousha, JA-C391-ippoutsuukou, JA-C391-tachidomaru, JA-C391-saegiru

\chapter{Crossing Safely}
\label{ch:ja-crossing-safely}
\hlchaptermodality{391}

\begin{chapteropening}
\textbf{By the end of this chapter:} \emph{I can cross the road safely.}

Complete the last lesson of chapter 391: I can cross the road safely.
\end{chapteropening}

\section[\texorpdfstring{\ja{おうだんほどう} (ōdanhodō)}{おうだんほどう (ōdanhodō)}]{\ja{おうだんほどう} (ōdanhodō) — a pedestrian crossing}

\label{lesson:JA-C391-oudanhodou}



\begin{quote}
Before the new one: say the Japanese for a licence, then the Japanese for to overtake.
\end{quote}

\subsection*{You'll want to know: \ja{おうだんほどう}}
\textbf{\ja{おうだんほどう}} — \emph{ōdanhodō} — \textquotedblleft{}a pedestrian crossing\textquotedblright{}.

\begin{quote}
\textbf{\ja{おうだんほどうを} \ja{わたって} \ja{ください。}} — \emph{ōdanhodō o watatte kudasai} — \textquotedblleft{}Please cross at the pedestrian crossing.\textquotedblright{}
\end{quote}

\begin{cousinweb}
\textbf{\ja{おうだんほどう}} (\emph{ōdanhodō}) is Sino-Japanese (kango): four Chinese characters read with their on'yomi, readings borrowed from Chinese. They are \emph{ō}, \textquotedblleft{}sideways\textquotedblright{} (Middle Chinese \emph{hwaeng}), \emph{dan}, \textquotedblleft{}severance\textquotedblright{} (\emph{dwanX}), \emph{ho}, \textquotedblleft{}walk\textquotedblright{} (\emph{buH}) and \emph{dō}, \textquotedblleft{}street\textquotedblright{} (\emph{dawX}) (KANJIDIC2; Middle Chinese in Baxter's transcription).
\end{cousinweb}

\begin{grammarlens}[title={What you've built}]
One more word for getting around.
\end{grammarlens}

\subsection*{Your turn}
\begin{itemize}
\raggedright
  \item \emph{Say it:} \emph{ōdanhodō}
  \item \emph{Say it:} \emph{ōdanhodō}, once more
  \item \emph{Say it:} \emph{ōdanhodō o watatte kudasai} — Please cross at the pedestrian crossing.
  \item \emph{Recall:} say the Japanese for a licence, then the Japanese for to overtake, then say \emph{ōdanhodō} again
\end{itemize}

\begin{tcolorbox}[breakable,colback=teal!4,colframe=teal!35!black,title={Before you move on}]
What does \ja{おうだんほどう} mean? (\textquotedblleft{}A pedestrian crossing\textquotedblright{}.) Say it once more.
\end{tcolorbox}
\section[\texorpdfstring{\ja{ほこうしゃ} (hokōsha)}{ほこうしゃ (hokōsha)}]{\ja{ほこうしゃ} (hokōsha) — a pedestrian}

\label{lesson:JA-C391-hokousha}



\begin{quote}
Before the new one: say the Japanese for to overtake, then the Japanese for a pedestrian crossing.
\end{quote}

\subsection*{You'll want to know: \ja{ほこうしゃ}}
\textbf{\ja{ほこうしゃ}} — \emph{hokōsha} — \textquotedblleft{}a pedestrian\textquotedblright{}.

\begin{quote}
\textbf{\ja{ほこうしゃは} \ja{ほどうを} \ja{あるきます。}} — \emph{hokōsha wa hodō o arukimasu} — \textquotedblleft{}Pedestrians walk on the pavement.\textquotedblright{}
\end{quote}

\begin{cousinweb}
\textbf{\ja{ほこうしゃ}} (\emph{hokōsha}) is Sino-Japanese (kango): three Chinese characters read with their on'yomi, readings borrowed from Chinese. They are \emph{ho}, \textquotedblleft{}walk\textquotedblright{} (Middle Chinese \emph{buH}), \emph{kō}, \textquotedblleft{}going\textquotedblright{} (\emph{haeng}) and \emph{sha}, \textquotedblleft{}someone\textquotedblright{} (\emph{tsyaeX}) (KANJIDIC2; Middle Chinese in Baxter's transcription).
\end{cousinweb}

\begin{grammarlens}[title={What you've built}]
One more word for getting around.
\end{grammarlens}

\subsection*{Your turn}
\begin{itemize}
\raggedright
  \item \emph{Say it:} \emph{hokōsha}
  \item \emph{Say it:} \emph{hokōsha}, once more
  \item \emph{Say it:} \emph{hokōsha wa hodō o arukimasu} — Pedestrians walk on the pavement.
  \item \emph{Recall:} say the Japanese for to overtake, then the Japanese for a pedestrian crossing, then say \emph{hokōsha} again
\end{itemize}

\begin{tcolorbox}[breakable,colback=teal!4,colframe=teal!35!black,title={Before you move on}]
What does \ja{ほこうしゃ} mean? (\textquotedblleft{}A pedestrian\textquotedblright{}.) Say it once more.
\end{tcolorbox}
\section[\texorpdfstring{\ja{いっぽうつうこう} (ippōtsūkō)}{いっぽうつうこう (ippōtsūkō)}]{\ja{いっぽうつうこう} (ippōtsūkō) — one-way traffic}

\label{lesson:JA-C391-ippoutsuukou}



\begin{quote}
Before the new one: say the Japanese for a pedestrian crossing, then the Japanese for a pedestrian.
\end{quote}

\subsection*{You'll want to know: \ja{いっぽうつうこう}}
\textbf{\ja{いっぽうつうこう}} — \emph{ippōtsūkō} — \textquotedblleft{}one-way traffic\textquotedblright{}.

\begin{quote}
\textbf{\ja{この} \ja{みちは} \ja{いっぽうつうこうです。}} — \emph{kono michi wa ippōtsūkō desu} — \textquotedblleft{}This street is one-way.\textquotedblright{}
\end{quote}

\begin{cousinweb}
\textbf{\ja{いっぽうつうこう}} (\emph{ippōtsūkō}) is Sino-Japanese (kango): four Chinese characters read with their on'yomi, readings borrowed from Chinese. They are \emph{ichi}, \textquotedblleft{}one\textquotedblright{} (Middle Chinese \emph{'jit}), \emph{hō}, \textquotedblleft{}direction\textquotedblright{} (\emph{pjang}), \emph{tsū}, \textquotedblleft{}traffic\textquotedblright{} (\emph{thuwng}) and \emph{kō}, \textquotedblleft{}going\textquotedblright{} (\emph{haeng}), joined as \emph{ippōtsūkō} (KANJIDIC2; Middle Chinese in Baxter's transcription).
\end{cousinweb}

\begin{grammarlens}[title={What you've built}]
One more word for getting around.
\end{grammarlens}

\subsection*{Your turn}
\begin{itemize}
\raggedright
  \item \emph{Say it:} \emph{ippōtsūkō}
  \item \emph{Say it:} \emph{ippōtsūkō}, once more
  \item \emph{Say it:} \emph{kono michi wa ippōtsūkō desu} — This street is one-way.
  \item \emph{Recall:} say the Japanese for a pedestrian crossing, then the Japanese for a pedestrian, then say \emph{ippōtsūkō} again
\end{itemize}

\begin{tcolorbox}[breakable,colback=teal!4,colframe=teal!35!black,title={Before you move on}]
What does \ja{いっぽうつうこう} mean? (\textquotedblleft{}One-way traffic\textquotedblright{}.) Say it once more.
\end{tcolorbox}
\section[\texorpdfstring{\ja{たちどまる} (tachidomaru)}{たちどまる (tachidomaru)}]{\ja{たちどまる} (tachidomaru) — to stop walking; to stand still}

\label{lesson:JA-C391-tachidomaru}



\begin{quote}
Before the new one: say the Japanese for a pedestrian, then the Japanese for one-way traffic.
\end{quote}

\subsection*{You'll want to know: \ja{たちどまる}}
\textbf{\ja{たちどまる}} — \emph{tachidomaru} — \textquotedblleft{}to stop walking; to stand still\textquotedblright{}.

\begin{quote}
\textbf{\ja{しんごうの} \ja{まえで} \ja{たちどまりました。}} — \emph{shingō no mae de tachidomarimashita} — \textquotedblleft{}I stopped in front of the traffic light.\textquotedblright{}
\end{quote}

Polite \textbf{\ja{たちどまります}}, \emph{tachidomarimasu}; a godan verb. Past (polite): \textbf{\ja{ほこうしゃが} \ja{たちどまりました。}}, \emph{hokōsha ga tachidomarimashita}.

\begin{cousinweb}
\textbf{\ja{たちどまる}} (\emph{tachidomaru}) is a native Japanese word (wago), inherited rather than borrowed from Chinese, built from \emph{tachi}, \textquotedblleft{}standing\textquotedblright{} (from \emph{tatsu}) and \emph{domaru}, \textquotedblleft{}stopping\textquotedblright{} (from \emph{tomaru}); its kanji were attached later as kun'yomi (JMdict).
\end{cousinweb}

\begin{grammarlens}[title={What you've built}]
One more word for getting around.
\end{grammarlens}

\subsection*{Your turn}
\begin{itemize}
\raggedright
  \item \emph{Say it:} \emph{tachidomaru}
  \item \emph{Say it:} \emph{tachidomaru}, once more
  \item \emph{Say it:} \emph{shingō no mae de tachidomarimashita} — I stopped in front of the traffic light.
  \item \emph{Recall:} say the Japanese for a pedestrian, then the Japanese for one-way traffic, then say \emph{tachidomaru} again
\end{itemize}

\begin{tcolorbox}[breakable,colback=teal!4,colframe=teal!35!black,title={Before you move on}]
What does \ja{たちどまる} mean? (\textquotedblleft{}To stop walking; to stand still\textquotedblright{}.) Say it once more.
\end{tcolorbox}
\section[\texorpdfstring{\ja{さえぎる} (saegiru)}{さえぎる (saegiru)}]{\ja{さえぎる} (saegiru) — to block; to obstruct}

\label{lesson:JA-C391-saegiru}



\begin{quote}
Before the new one: say the Japanese for one-way traffic, then the Japanese for to stop walking.
\end{quote}

\subsection*{You'll want to know: \ja{さえぎる}}
\textbf{\ja{さえぎる}} — \emph{saegiru} — \textquotedblleft{}to block; to obstruct\textquotedblright{}.

\begin{quote}
\textbf{\ja{くるまが} \ja{みちを} \ja{さえぎっていました。}} — \emph{kuruma ga michi o saegitte imashita} — \textquotedblleft{}A car was blocking the road.\textquotedblright{}
\end{quote}

Polite \textbf{\ja{さえぎります}}, \emph{saegirimasu}; a godan verb. Past (polite): \textbf{\ja{くるまが} \ja{どうろを} \ja{さえぎりました。}}, \emph{kuruma ga dōro o saegirimashita}.

\begin{cousinweb}
\textbf{\ja{さえぎる}} (\emph{saegiru}) is a native Japanese word (wago), inherited rather than borrowed from Chinese; its kanji was attached later as a kun'yomi (JMdict).
\end{cousinweb}

\begin{grammarlens}[title={What you've built}]
One more word for getting around.
\end{grammarlens}

\subsection*{Your turn}
\begin{itemize}
\raggedright
  \item \emph{Say it:} \emph{saegiru}
  \item \emph{Say it:} \emph{saegiru}, once more
  \item \emph{Say it:} \emph{kuruma ga michi o saegitte imashita} — A car was blocking the road.
  \item \emph{Recall:} say the Japanese for one-way traffic, then the Japanese for to stop walking, then say \emph{saegiru} again
\end{itemize}

\begin{tcolorbox}[breakable,colback=teal!4,colframe=teal!35!black,title={Before you move on}]
What does \ja{さえぎる} mean? (\textquotedblleft{}To block; to obstruct\textquotedblright{}.) Say it once more.
\end{tcolorbox}
